% from ega3-1-fr.tex
\oldpage[III]{100}

\begin{theorem}[2.2.1]
\label{III.2.2.1-fr}
(Serre). Soient $Y$ un préschéma noethérien, $f:X\to Y$ un morphisme
propre, $\mathcal L$ un $\mathcal O_X$-Module inversible ample pour $f$.
Pour tout $\mathcal O_X$-Module $\mathcal F$, posons
\[
\mathcal F(n)=\mathcal F\otimes_{\mathcal O_X}\mathcal L^{\otimes n}
\qquad\text{pour tout }n\in\mathbb Z.
\]
Alors, pour tout $\mathcal O_X$-Module cohérent $\mathcal F$~:
\begin{enumerate}
\item[{\rm(i)}] Les $R^qf_*(\mathcal F)$ sont des
$\mathcal O_Y$-Modules cohérents.
\item[{\rm(ii)}] Il existe un entier $N$ tel que pour $n\geq N$, on ait
\[
R^qf_*(\mathcal F(n))=0\qquad\text{pour tout }q>0.
\]
\item[{\rm(iii)}] Il existe un entier $N$ tel que, pour $n\geq N$,
l'homomorphisme canonique
\[
f^*(f_*(\mathcal F(n)))\longrightarrow\mathcal F(n)
\]
soit surjectif.
\end{enumerate}
\end{theorem}

\begin{proof}
Notons d'abord que si le théorème est vrai quand on y remplace
$\mathcal L$ par $\mathcal L^{\otimes d}$ $(d>0)$, il est vrai sous sa
forme initiale. En effet, on peut alors écrire
\[
\mathcal F(n)
=(\mathcal F\otimes\mathcal L^{\otimes r})
\otimes\mathcal L^{\otimes hd}
\]
avec un $h>0$ et $0\leq r<d$, et par hypothèse pour chaque $r$ il y a
un entier $N_r$ tel que pour $h\geq N_r$ les propriétés (ii) et (iii)
aient lieu pour le $\mathcal O_X$-Module
$\mathcal F\otimes\mathcal L^{\otimes r}$~; prenant pour $N$ le plus grand
des $dN_r$, (ii) et (iii) auront lieu pour $n\geq N$. On peut donc
supposer $\mathcal L$ très ample relativement à $f$ (II, 4.6.11)~; il
existe par suite une $Y$-immersion ouverte dominante
$i:X\to P$, où $P=\operatorname{Proj}(\mathcal S)$, où $\mathcal S$ est
une $\mathcal O_Y$-Algèbre quasi-cohérente graduée à degrés positifs, dans
laquelle $\mathcal S_1$ est de type fini et engendre $\mathcal S$~; en
outre, $\mathcal L$ est isomorphe à $i^*(\mathcal O_P(1))$ (II, 4.4.7).
Mais comme $f$ est propre, il en est de même de $i$ (II, 5.4.4), donc
$i$ est un isomorphisme $X\xrightarrow{\sim}P$. On peut donc se borner au
cas où $X=\operatorname{Proj}(\mathcal S)$ et
$\mathcal L=\mathcal O_X(1)$. Le th. (2.2.1) est alors conséquence de la
\end{proof}

\begin{proposition}[2.2.2]
\label{III.2.2.2-fr}
Soient $A$ un anneau noethérien, $S$ une $A$-algèbre graduée à degrés
positifs, dans laquelle $S_1$ est un $A$-module ayant un système de $r+1$
générateurs, et qui engendre l'algèbre $S$. Soit
$X=\operatorname{Proj}(S)$. Pour tout $\mathcal O_X$-Module cohérent
$\mathcal F$~:
\begin{enumerate}
\item[{\rm(i)}] Les $A$-modules $H^q(X,\mathcal F)$ sont de type fini.
\item[{\rm(ii)}] On a $H^q(X,\mathcal F)=0$ pour $q>r$.
\item[{\rm(iii)}] Il existe un entier $N$ tel que pour $n\geq N$, on ait
$H^q(X,\mathcal F(n))=0$ pour tout $q>0$.
\item[{\rm(iv)}] Il existe un entier $N$ tel que pour $n\geq N$,
$\mathcal F(n)$ soit engendré par ses sections au-dessus de $X$.
\end{enumerate}
\end{proposition}

\begin{proof}
Montrons d'abord comment (2.2.2) entraîne (2.2.1)~: dans (2.2.1)
(ramené au cas particulier $X=\operatorname{Proj}(\mathcal S)$ considéré
ci-dessus), $Y$ est quasi-compact, donc peut être recouvert par un nombre
fini d'ouverts affines, d'anneaux noethériens, tels que la restriction de
$\mathcal S_1$ à chacun de ces ouverts $U_\alpha$ soit engendrée par un
nombre fini de sections de $\mathcal S_1$ au-dessus de $U_\alpha$. Si on
suppose (2.2.2) démontré, il suffira alors de prendre pour $N$ dans les
parties (ii) et (iii) de (2.2.1) le plus grand des entiers analogues
correspondant aux $U_\alpha$ (tenant compte de (1.4.11) et de
(II, 3.4.7)).

Pour prouver (2.2.2), remarquons que $X$ s'identifie à un sous-schéma
fermé de $P=\mathbf P_A^r$ (II, 3.6.2)~; en outre, si $j:X\to P$ est
l'injection canonique, $j_*(\mathcal F)$ est un
$\mathcal O_P$-Module cohérent et on a
$j_*(\mathcal F(n))=(j_*(\mathcal F))(n)$ (II, 3.4.5 et 3.5.2).
Compte tenu de (G, II, cor. du th. 4.9.1), on est donc ramené à prouver
(2.2.2) dans le cas particulier où
\[
X=\mathbf P_A^r\quad\text{et}\quad S=A[T_0,\ldots,T_r].
\]

